%% kapa: exam=1 round=23 year=2012 period=1 form=A subject=civil_law subject_ko=민법 seq=39 number=39
%% source: Q-Net 감정평가사 기출문제 (cst00309) artlSeq=5052578 file=raw/1008/5052578/2012년 제23회 감정평가사 제1차시험 기출문제(1교시 A형).pdf
%% conversion: kordoc 4.24.51 -> tools/gen1.py (auto); review=proofread-sonnet; flags=none
\begin{question}{39}
부동산 소유권 이전등기의 추정력에 관한 설명으로 옳지 않은 것은? (다툼이 있으면 판례에 의함)
\begin{choices}
  \item 등기명의자는 그 전 소유자에 대하여는 적법한 등기원인에 의하여 소유권을 취득한 것이라는 등기의 추정력을 주장할 수 없다.
  \item 등기명의자가 등기부와 다른 등기원인을 주장하였으나 그 주장 사실이 인정되지 않은 것만으로는 등기의 추정력이 깨어지는 것은 아니다.
  \item 소유권이전등기의 원인으로 주장된 계약서가 진정하지 않은 것으로 증명된 경우, 그 등기의 적법 추정은 깨어진다.
  \item 등기원인 사실에 대한 입증이 부족하다는 이유로 그 등기를 무효라고 단정할 수 없다.
  \item 증여를 원인으로 미성년자에게서 친권자 앞으로 소유권이전등기가 경료된 경우, 특별한 사정이 없는 한 그 이전등기에 관하여 필요한 절차를 적법하게 거친 것으로 추정된다.
\end{choices}
\end{question}
\answer{1}
